\section{Extremal partitions and their centroids}
\label{ext:section}

Our goal is to replace an arbitrary measurable partition by an optimal
partition described by finitely many Gaussian linear scores. We include
the extremal reduction of Khot and
Naor~\cite[Theorem~1.2 and Lemma~3.3]{KhotNaor} in the form needed below.

It suffices to prove the upper bound for $n+1$ cells in
$\mathbb R^n$ with $n\ge4$. Indeed, given $k$ cells in
$\mathbb R^d$, put $n=\max\{d,k-1,4\}$, take their products with
$\mathbb R^{n-d}$, and append $n+1-k$ empty cells. Gaussian product
decomposition appends zero coordinates to each first moment and
preserves the objective.

Fix such an $n$ and put $N=n+1$. All partitions below are measurable
partitions up to Gaussian null sets, and empty cells are permitted. Write
\[
 C=\sup_{(A_1,\ldots,A_N)}
       \sum_{i=1}^N\left\|\int_{A_i}x\,d\gamma_n(x)\right\|^2,
 \qquad B=\frac{3}{2\sqrt{2\pi}}.
\]
The moments exist, and Cauchy--Schwarz on each cell gives $C\le n$.
We first record the normalized vector form of the Gaussian-maximum
description in \cite[Lemma~3.7]{KhotNaor}. Its proof also gives
attainment directly.

\begin{lemma}\label{ext:duality}
If $g$ is a standard Gaussian vector in $\mathbb R^n$, then
\begin{equation}\label{ext:dual-formula}
 \sqrt C=
 \max_{\sum_{i=1}^N\|v_i\|^2=1}
       \mathbb E\max_{1\le i\le N}\langle v_i,g\rangle.
\end{equation}
The supremum defining $C$ is attained. Moreover, for every $1\le k\le N$
and every list $v_1,\ldots,v_k\in\mathbb R^n$,
\begin{equation}\label{ext:shorter-list}
 \mathbb E\max_{1\le i\le k}\langle v_i,g\rangle
 \le \sqrt C\left(\sum_{i=1}^k\|v_i\|^2\right)^{1/2}.
\end{equation}
\end{lemma}

\begin{proof}
For a partition with moments $z_i=\int_{A_i}x\,d\gamma_n(x)$,
Euclidean norm duality gives
\[
 \left(\sum_i\|z_i\|^2\right)^{1/2}
 =\max_{\sum_i\|v_i\|^2=1}
       \sum_i\int_{A_i}\langle v_i,x\rangle\,d\gamma_n(x).
\]
Taking the supremum jointly over partitions and lists, and then first
over partitions, yields the right side of \eqref{ext:dual-formula}.
Indeed, for each fixed list the pointwise maximum is attained by assigning
$x$ to a largest score $\langle v_i,x\rangle$, with ties resolved by the
smallest index. This gives a measurable partition. The expected maximum
is continuous in the list: the absolute difference of two maxima is at
most $\|g\|\max_i\|v_i-w_i\|$. Compactness of the sphere of lists therefore
gives a maximizing list. Its score partition has moment norm at least
its pairing with that list, namely $\sqrt C$, so it attains $C$.

For \eqref{ext:shorter-list}, let $Q_1,\ldots,Q_k$ be a score partition
for the specified list, and put $b_i=\int_{Q_i}x\,d\gamma_n(x)$.
Appending $N-k$ empty cells shows that $\sum_i\|b_i\|^2\le C$.
Consequently
\[
 \mathbb E\max_i\langle v_i,g\rangle
 =\sum_i\langle v_i,b_i\rangle
 \le\left(\sum_i\|v_i\|^2\right)^{1/2}
      \left(\sum_i\|b_i\|^2\right)^{1/2},
\]
which proves the assertion.
\end{proof}

We argue by contradiction and henceforth assume
\begin{equation}\label{ext:contradiction}
 C>B^2=\frac{9}{8\pi}.
\end{equation}
Choose an optimal partition with the fewest cells of positive measure,
discard its null cells, and call the remaining cells \emph{active}.
Denote their moments by
$z_1,\ldots,z_m$. Thus
\begin{equation}\label{ext:centroid-identities}
 \sum_{i=1}^m z_i=0,
 \qquad \sum_{i=1}^m\|z_i\|^2=C.
\end{equation}

The following structure statement develops the minimal-cell reduction
of \cite[Lemma~3.3]{KhotNaor}; the rank assertions are proved directly
from the Gram matrix.

\begin{lemma}\label{ext:structure}
The centroids satisfy $\langle z_i,z_j\rangle<0$ for $i\ne j$.
Their span $V$ has dimension $m-1$, every proper subset of the centroids
is linearly independent, and, for each fixed $i$, the $m-1$ vectors
$z_i-z_j$, $j\ne i$, form a basis of $V$.
The optimal cells have the closed conical representatives
\begin{equation}\label{ext:cones}
 K_i=\{x\in\mathbb R^n:
          \langle z_i-z_j,x\rangle\ge0\text{ for every }j\ne i\}.
\end{equation}
Furthermore, $m\ge5$.
\end{lemma}

\begin{proof}
Since $C>0$, we have $m\ge2$. Merging two cells changes the objective by
$2\langle z_i,z_j\rangle$. A positive product would contradict optimality;
a zero product would contradict minimality of $m$. Thus every such
product is negative, and in particular every $z_i$ is nonzero.

For the original optimal partition, pairing each cell with its own
centroid gives
\[
 C\le\mathbb E\max_i\langle z_i,g\rangle\le C,
\]
where the second inequality is \eqref{ext:shorter-list}.
The difference between the maximum score and the assigned score is
nonnegative and has zero integral, so it vanishes almost everywhere.
The centroids are distinct, and their score ties lie on finitely many
Gaussian-null hyperplanes. This proves \eqref{ext:cones}; overlaps
between these closed representatives are immaterial.

Let $Q=(\langle z_i,z_j\rangle)_{i,j=1}^m$ be the Gram matrix.
Its row sums vanish by \eqref{ext:centroid-identities}. Hence, for every
$a=(a_1,\ldots,a_m)\in\mathbb R^m$,
\begin{equation}\label{ext:laplacian}
 a^{\mathsf T}Qa
 =\sum_{i<j}(-\langle z_i,z_j\rangle)(a_i-a_j)^2.
\end{equation}
All weights are positive, so the kernel of $Q$ consists exactly of
constant vectors. Its rank is $m-1$, and the only linear relations among
the centroids are multiples of their sum. A dependence on a proper
subset would have a zero coefficient and therefore must be trivial.
A relation among $z_i-z_j$, $j\ne i$, gives a relation among all centroids
whose coefficients sum to zero; it too must be trivial. This proves
the linear-algebra assertions.

Each $K_i$ depends only on the orthogonal projection onto $V$.
If $d=\dim V$, then orthogonal Gaussian product decomposition identifies
$K_i$ with $(K_i\cap V)\times V^\perp$, and its moment is the moment of
$K_i\cap V$ in $V$, extended by zero on $V^\perp$.
If $m\le4$, then $d\le3$. Taking the product of this partition of
$V\cong\mathbb R^d$ with $\mathbb R^{3-d}$ preserves its moment norms.
The three-dimensional propeller theorem of Heilman, Jagannath, and
Naor~\cite[Theorem~1.1]{HJN}, which permits any finite number of cells,
would then give $C\le9/(8\pi)$, contradicting
\eqref{ext:contradiction}. Thus $m\ge5$.
\end{proof}

We have reduced a hypothetical counterexample to at least five active
cells whose centroid span has dimension one less than their number.
From now on we work in $V\cong\mathbb R^{m-1}$, use $K_i$ for the reduced
cones, and let $g$ be standard Gaussian in this space. Their moments
remain exactly $z_i$. Inequality~\eqref{ext:shorter-list} remains valid
for lists in $V$, by embedding them in $\mathbb R^n$. Set
\begin{equation}\label{ext:notation}
 X_i=\langle z_i,g\rangle,
 \qquad P_i=\gamma_{m-1}(K_i),
 \qquad r_i=\frac{\|z_i\|}{\sqrt C}.
\end{equation}
Then $P_i>0$, $r_i>0$, $\sum_iP_i=\sum_i r_i^2=1$, and
$\mathbb E\max_iX_i=C$. Since $\sum_iX_i=0$, a winning score is
nonnegative. Thus $K_i\subseteq\{X_i\ge0\}$ and
\begin{equation}\label{ext:half-probability}
 P_i\le\frac12\qquad(1\le i\le m).
\end{equation}
